\documentclass[12pt]{article}
\usepackage[utf8]{inputenc}
\usepackage[T1]{fontenc}
\usepackage{hyperref}
\title{Upper Saxon Language Resources for K-12 Curriculum Development}
\author{Compiled by Claude Sonnet 5.5 (claude-sonnet-5-5) for Andr\'as Kornai}
\date{}
\begin{document}
\maketitle

\section{Introduction}
Upper Saxon (Obers\"achsisch; ISO 639-3 sxu) is an East Central German variety spoken in Saxony and neighbouring areas of Germany. This report is based on a single web search. Little dedicated teaching material turned up, and nothing found is a school curriculum in Upper Saxon. Wikipedia describes the traditional dialect as now largely replaced by a regional colloquial form (regiolect).

\section{Primary Educational Resources}
No primer, textbook or school programme for Upper Saxon was found in the search results.

\section{Linguistic and Reference Materials}
\begin{itemize}
\item Upper Saxon German, Wikipedia overview (English): \url{https://en.wikipedia.org/wiki/Upper_Saxon_German}
\item Area: Obers\"achsisch, phonetic data page of the regionalsprache.de PhonD2 database, with a recorded dialect sample: \url{https://regionalsprache.de/phonD2/138212.html}
\item \emph{S\"achsisch}, a humorous dialect dictionary (Duden), sample PDF: \url{https://shop.duden.de/media/44/24/e4/1719879356/Leseprobe_9783411756841_Sächsisch.pdf}
\end{itemize}
Caution: the \emph{Altsächsisches Handwörterbuch} that appears in search results covers Old Saxon, a Low German language, and is not Upper Saxon.

\section{Religious and Community Resources}
None were found in the search results.

\section{Digital and Audio Resources}
The only audio found is the recorded dialect sample on the PhonD2 page listed above.

\section{Curriculum Development Context}
Standard German is the school language. Search results gave no evidence of Upper Saxon being taught in schools. A curriculum would have to be built from dialect documentation (such as the regional atlases and recordings above), which would need further searching and consultation with speakers.
\end{document}